
\subsubsection{Parameter-Efficient Fine-Tuning}

Low-Rank Adaptation (LoRA) introduced the decomposition $\Delta W$ = BA to constrain weight updates to a low-rank subspace. Subsequent work has focused on improving LoRA training: DoRA decomposes weights into magnitude and direction components; QLoRA enables efficient fine-tuning of quantized models; AdaLoRA adaptively allocates rank across layers. StelLA employs Stiefel manifold geometry to optimize LoRA training. These methods leverage geometric structure to improve training but do not analyze trained adapters for task similarity.

\subsubsection{Weight-Space Learning}

Weight-space learning treats neural network parameters as data points for downstream analysis. The Model Zoo framework introduced controlled populations of neural networks with verified provenance, demonstrating that weight-space representations enable model property prediction, generation, and analysis. This methodology has focused primarily on full neural network models rather than parameter-efficient adapters. The present work adapts this controlled-generation approach specifically for LoRA adapters.

\subsubsection{Geometric Analysis of Representations}

Grassmann manifolds provide a natural framework for comparing linear subspaces. Principal angles between subspaces, computed via SVD of inner products between orthonormal bases, yield the Grassmann geodesic distance. This metric has been applied to subspace clustering, domain adaptation, and representation similarity analysis. CKA and SVCCA measure representation similarity across layers and models by comparing activations rather than weights. For LoRA adapters, weight-space analysis requires no inference: geometric signatures can be computed directly from adapter parameters.

FL-TAC clusters similar adapters in federated learning settings, though without controlled provenance verification. The present work provides controlled validation where only the training task varies.
